\documentclass[border=0pt]{standalone}
% The shared preamble of every docs diagram (docs/diagrams/<name>.tex).
% dvisvgm's driver first, so patterns and opacity survive into the SVG.
\def\pgfsysdriver{pgfsys-dvisvgm.def}
\usepackage[siunitx,RPvoltages]{circuitikz}
\usepackage{xcolor}
% Latin Modern: outline fonts for every glyph a label uses, µ included.
\usepackage[T1]{fontenc}
\usepackage{lmodern}
\usetikzlibrary{arrows.meta,backgrounds,calc,patterns,positioning}
\tikzset{
  % A white card, so diagrams read in mdBook's dark themes too.
  every picture/.append style={show background rectangle,
    background rectangle/.style={fill=white,inner sep=4pt}},
  % Layout geometry: layers told apart by fill and hatching, never colour alone.
  layer a/.style={draw=black,line width=0.6pt,fill=gray!20},
  layer b/.style={draw=black,line width=0.6pt,fill=gray!45},
  layer c/.style={draw=black,line width=0.6pt,fill=white,postaction={pattern=north east lines}},
  layer d/.style={draw=black,line width=0.6pt,fill=white,postaction={pattern=crosshatch dots}},
  marker/.style={draw=black,dashed,line width=0.6pt},
  % A measured distance and its label.
  dim/.style={{Latex[length=2mm]}-{Latex[length=2mm]},line width=0.5pt},
  dim label/.style={font=\small,fill=white,inner sep=1pt},
  % Where a rule fails: thick, dashed and labelled, not only red.
  violation/.style={draw=red!75!black,line width=1.4pt,dash pattern=on 3pt off 2pt},
  violation label/.style={font=\small\bfseries,text=red!75!black},
  caption/.style={font=\small},
}

\begin{document}
\begin{circuitikz}[font=\small]
  % Left: the reference card, a regular nfet. Right: the drawn transistor
  % with the hvtp implant over its gate: ngate and hvtp, so nfet_hvt.
  \node[font=\bfseries] at (1.2,3.6) {Reference netlist};
  \node[nmos] (n) at (1.0,1.6) {};
  \draw (n.G) -- ++(-0.3,0) node[left] {G};
  \draw (n.D) -- ++(0,0.3) node[above] {D};
  \draw (n.S) -- ++(0,-0.3) node[below] {S};
  \node[right,violation] at (1.5,1.9) {M1 nfet};
  \node[right] at (1.5,1.3) {W \SI{1}{\micro\metre}, L \SI{0.15}{\micro\metre}};

  \node[font=\bfseries] at (7.0,3.6) {Layout, extracted};
  \draw[layer a] (5.4,0.6) rectangle (8.6,2.6);
  \draw[layer d] (6.5,0.1) rectangle (7.5,3.1);
  \draw[layer c] (6.75,0.2) rectangle (7.25,3.0);
  \draw[marker] (6.75,0.6) rectangle (7.25,2.6);
  \node[font=\footnotesize] at (6.0,1.6) {S};
  \node[font=\footnotesize] at (8.0,1.6) {D};
  \node[font=\footnotesize,right] at (7.5,2.9) {hvtp};
  \node[violation label,align=center] at (7.0,-0.45) {\texttt{ngate and hvtp}: model nfet\_hvt\\ against nfet};
  \node[caption] at (4.6,-1.4) {grey: diffusion; hatched: poly; dotted: the hvtp implant; dashed: the gate marker};
\end{circuitikz}
\end{document}
